\section{Exact flux spectrum on a cycle}
\label{sec:flux}

Let \(P\) be the adjacency matrix of the unweighted cycle \(C_n\),
\(n\ge3\), so \(\rho(P)=2\).  Let \(|H|=P\), and suppose the total
oriented holonomy is \(e^{\mathrm i\Theta}\).

\begin{theorem}[Flux-cycle spectrum]\label{thm:flux}
After switching and uniformly distributing the flux,
\begin{equation}\label{eq:flux-spectrum}
 \spec(H)
 =
 \left\{
 2\cos\left(\frac{\Theta+2\pi k}{n}\right):
 0\le k<n
 \right\}.
\end{equation}
Define
\[
 d_+(\Theta)=\dist(\Theta,2\pi\mathbb Z),
 \qquad
 d_-(\Theta)=\dist(\Theta,n\pi+2\pi\mathbb Z),
\]
with distances represented in \([0,\pi]\).  Then
\begin{align}
 \delta_+
 &=
 2\left(1-\cos\frac{d_+}{n}\right),
 \label{eq:flux-plus}\\
 \delta_-
 &=
 2\left(1-\cos\frac{d_-}{n}\right),
 \label{eq:flux-minus}\\
 2-\|H\|
 &=
 2\min_{\sigma\in\{+,-\}}
 \left(1-\cos\frac{d_\sigma}{n}\right).
 \label{eq:flux-norm}
\end{align}
\end{theorem}

\begin{proof}
Switch all but one edge to phase \(1\), then conjugate by a diagonal
unitary that distributes the remaining phase equally around the cycle.
The Fourier vectors diagonalize the resulting twisted circulant and
give \eqref{eq:flux-spectrum}.  Maximizing the cosine gives
\eqref{eq:flux-plus}.  Multiplication of every edge by \(-1\) changes
total holonomy by \((-1)^n=e^{\mathrm i n\pi}\), which gives
\eqref{eq:flux-minus}.  Equation \eqref{eq:norm-defect} gives the last
formula.
\end{proof}

\begin{corollary}[Quadratic small-flux law]\label{cor:quadratic}
For either sector,
\begin{equation}\label{eq:quadratic}
 \frac{4d_\sigma^2}{\pi^2n^2}
 \le\delta_\sigma
 \le\frac{d_\sigma^2}{n^2}.
\end{equation}
\end{corollary}

\begin{proof}
Use
\[
 \frac{2t^2}{\pi^2}\le1-\cos t\le\frac{t^2}{2}
 \qquad(|t|\le\pi).
\]
\end{proof}

Thus a nontrivial but shrinking flux can make every finite matrix
strict while producing no uniform power-scale gap.  For an even cycle,
both sectors are balanced exactly at holonomy \(1\).  For an odd cycle,
operator-norm equality can occur at holonomy \(1\) through the positive
sector or at holonomy \(-1\) through the negative sector.
